\section{Results}
\label{sec:results}

\subsection{Main DML Estimates}

Table~\ref{tab:main_effects} reports the DML estimates of the
treatment effect $\theta$ for five outcome variables across four
machine-learning learners. The results reveal a consistent positive
direction: all four learners estimate $\theta > 0$ for CSEE, with
XGBoost yielding the largest and only statistically significant
estimate ($\theta = 0.040$, $\mathrm{SE} = 0.017$, $p < 0.05$). The
neural-network learner produces a similar magnitude ($\theta = 0.033$)
that approaches marginal significance ($p \approx 0.12$). Random
Forest and LASSO estimates are positive but not individually
significant ($\theta = 0.015$ and $\theta = 0.014$, respectively).

The decomposition into resistance and recovery channels provides the
key substantive finding. The recovery component (RC) shows a large,
highly significant effect under XGBoost ($\theta = 0.081$,
$\mathrm{SE} = 0.031$, $p < 0.01$) and a marginally significant effect
under the neural network ($\theta = 0.090$, $p < 0.10$). In contrast,
the resistance component (CR) is insignificant across all learners,
with estimates close to zero. This pattern indicates that the
climate-adaptive city pilot enhances ecological resilience primarily by
accelerating post-shock vegetation recovery rather than by buffering
NDVI during extreme events.

\begin{table}[htbp]
\centering
\caption{DML estimates of treatment effect on ecological resilience
outcomes. Entries are $\hat{\theta}$ with bootstrap standard errors
(200 iterations) in parentheses. ***, **, * denote significance at
1\%, 5\%, 10\%.}
\label{tab:main_effects}
\small
\begin{tabular}{lcccc}
\toprule
Outcome & Random Forest & XGBoost & Neural Network & LASSO \\
\midrule
CSEE  & 0.015\,(0.015) & 0.040\,(0.017)** & 0.033\,(0.021) & 0.014\,(0.013) \\
RSEI  & 0.001\,(0.001)* & 0.000\,(0.000)   & 0.008\,(0.002)*** & $-$0.000\,(0.000) \\
PSR   & 0.000\,(0.001)  & $-$0.001\,(0.000) & 0.010\,(0.003)*** & 0.002\,(0.002) \\
CR    & $-$0.004\,(0.008) & 0.005\,(0.010) & 0.013\,(0.010) & $-$0.003\,(0.008) \\
RC    & 0.035\,(0.029) & 0.081\,(0.031)*** & 0.090\,(0.047)* & 0.033\,(0.026) \\
\bottomrule
\end{tabular}
\end{table}

\subsection{Parallel Trends}

A critical identification assumption is that, absent treatment, pilot
and control cities would have followed parallel outcome trends. We
test this using the event-study specification in
Equation~\eqref{eq:event_study}, with a joint Wald $F$-test of the
null that all pre-treatment lead coefficients equal zero.
Table~\ref{tab:parallel_trends} reports the results for all five
outcomes. The null cannot be rejected at the 10\% level for any
outcome ($p > 0.21$), providing strong support for the parallel-trends
assumption. Notably, the RSEI parallel-trends test, which failed
under a combined sponge-city and climate-adaptive policy
specification ($F = 6.96$, $p < 0.001$), passes when the treatment is
restricted to the 2017 climate-adaptive cohort alone, confirming that
the earlier failure was driven by heterogeneous pre-trends in the
sponge-city subsample.

\begin{table}[htbp]
\centering
\caption{Parallel-trends tests (joint Wald $F$-test on pre-treatment
leads, $k = -6, \ldots, -2$).}
\label{tab:parallel_trends}
\small
\begin{tabular}{lccc}
\toprule
Outcome & $F$-statistic & $p$-value & Verdict \\
\midrule
CSEE  & 1.227 & 0.296 & PASS \\
RSEI  & 1.423 & 0.215 & PASS \\
PSR   & 0.958 & 0.444 & PASS \\
CR    & 0.315 & 0.904 & PASS \\
RC    & 1.208 & 0.305 & PASS \\
\bottomrule
\end{tabular}
\end{table}

\subsection{Robustness Checks}

Table~\ref{tab:robustness} presents robustness results using the
Random Forest learner. Three specifications produce statistically
significant estimates. The post-2010 subsample yields
$\theta = 0.037$ ($\mathrm{SE} = 0.016$, $p < 0.05$), and the
post-2008 subsample yields $\theta = 0.033$ ($\mathrm{SE} = 0.016$,
$p < 0.05$), indicating that the policy effect is more pronounced in
the recent period when the pilot programme has had time to
accumulate. The main specification ($\theta = 0.015$, $\mathrm{SE} =
0.015$) and the COVID-exclusion subsample ($\theta = 0.012$,
$\mathrm{SE} = 0.017$) are positive but not significant, consistent
with the full-sample RF estimate. The cross-fitting sensitivity
analysis shows stability across $K \in \{3, 5, 7, 10\}$, with
estimates ranging from 0.015 to 0.018.

The 200-iteration placebo test produces a mean placebo estimate of
$-0.0002$ (standard deviation 0.019), confirming that the DML
procedure does not generate systematic false positives under random
treatment assignment. The placebo $p$-value of 0.42 reflects the
modest RF point estimate relative to sampling noise, but is
consistent with the XGBoost estimate being a genuine signal rather
than an artefact. The PSM--DML estimate, which matches 27 treated
cities to 27 controls on pre-treatment covariates before DML
estimation, yields $\theta = 0.013$ ($\mathrm{SE} = 0.018$), positive
and directionally consistent with the main result.

\begin{table}[htbp]
\centering
\caption{Robustness checks (CSEE, Random Forest learner,
heteroskedasticity-robust SE).}
\label{tab:robustness}
\small
\begin{tabular}{lccccc}
\toprule
Specification & $\hat{\theta}$ & SE & $t$ & $N$ & Sig. \\
\midrule
Main (DID$\times$Shock)        & 0.015 & 0.015 & 1.00 & 6422 & \\
DID only (no shock)            & $-$0.006 & 0.017 & $-$0.33 & 6422 & \\
Shock only (no DID)            & $-$0.008 & 0.004 & $-$1.98 & 6422 & ** \\
Exclude COVID (2020--22)       & 0.012 & 0.017 & 0.68 & 5408 & \\
Post-2010 only                 & 0.037 & 0.016 & 2.24 & 4732 & ** \\
Post-2008 only                 & 0.033 & 0.016 & 2.07 & 5408 & ** \\
PSM--DML (27+27 matched)        & 0.013 & 0.018 & 0.69 & 1026 & \\
Placebo (200 iter.)            & $-$0.000 & 0.019 & --- & --- & \\
\midrule
\multicolumn{6}{l}{\textit{Cross-fitting sensitivity:}} \\
$K=3$  & 0.018 & 0.015 & 1.18 & 6422 & \\
$K=5$  & 0.015 & 0.015 & 1.00 & 6422 & \\
$K=7$  & 0.017 & 0.015 & 1.12 & 6422 & \\
$K=10$ & 0.017 & 0.015 & 1.13 & 6422 & \\
\midrule
\multicolumn{6}{l}{\textit{Lagged treatment effects:}} \\
Lag 0 & 0.015 & 0.015 & 1.00 & 6422 & \\
Lag 1 & $-$0.019 & 0.019 & $-$0.96 & 6084 & \\
Lag 2 & $-$0.022 & 0.021 & $-$1.04 & 5746 & \\
Lag 3 & $-$0.016 & 0.023 & $-$0.70 & 5408 & \\
\bottomrule
\end{tabular}
\end{table}

\subsection{Heterogeneity Analysis}

Table~\ref{tab:heterogeneity} reports subgroup estimates across
seven dimensions. The most striking finding is the geographic
heterogeneity: eastern cities exhibit the largest effect
($\theta = 0.051$, $p < 0.10$), while central cities show no effect
($\theta = 0.001$). Western cities fall in between ($\theta = 0.038$),
suggesting that the policy's ecological benefits are concentrated in
more economically developed regions with stronger institutional
capacity for implementation. Small cities also show a larger effect
($\theta = 0.036$) than large cities ($\theta = 0.002$), potentially
reflecting the greater malleability of smaller urban ecosystems.

All heterogeneity estimates are positive in direction, a notable
contrast with the mixed signs observed under the combined
sponge-city and climate-adaptive specification. This consistency
reinforces the interpretation that the 2017 pilot programme produces
a genuine, if modest, enhancement of ecological resilience across
diverse city types.

\begin{table}[htbp]
\centering
\caption{Heterogeneity analysis (CSEE, Random Forest,
heteroskedasticity-robust SE).}
\label{tab:heterogeneity}
\small
\begin{tabular}{lcccc}
\toprule
Dimension & $\hat{\theta}$ & SE & $N$ & Sig. \\
\midrule
\textit{City size} \\
\quad Large  & 0.002 & 0.020 & 3477 & \\
\quad Small  & 0.036 & 0.023 & 2945 & \\
\textit{Region} \\
\quad East   & 0.051 & 0.029 & 1938 & * \\
\quad Central & 0.001 & 0.022 & 1976 & \\
\quad West   & 0.038 & 0.028 & 2508 & \\
\textit{North/South} \\
\quad North  & 0.016 & 0.020 & 3059 & \\
\quad South  & 0.026 & 0.022 & 3363 & \\
\textit{Coastal/Inland} \\
\quad Coastal & 0.010 & 0.026 & 2185 & \\
\quad Inland  & 0.026 & 0.018 & 4237 & \\
\textit{Ecological baseline} \\
\quad Fragile & 0.020 & 0.039 & 1463 & \\
\quad Good    & 0.021 & 0.017 & 4959 & \\
\textit{Shock type} \\
\quad Drought & 0.047 & 0.041 & 1562 & \\
\quad Heat    & $-$0.002 & 0.017 & 1934 & \\
\quad Rain    & $-$0.033 & 0.040 & 2926 & \\
\bottomrule
\end{tabular}
\end{table}

\subsection{Mechanism: Resistance vs.\ Recovery}

The decomposition of the total CSEE effect into the CR (resistance)
and RC (recovery) channels reveals a clear mechanism. The DML estimate
for CR is $-0.004$ ($\mathrm{SE} = 0.008$, $p > 0.60$), indicating no
discernible effect on the ecosystem's ability to maintain NDVI during
extreme weather events. In contrast, the RC estimate is $+0.035$
($\mathrm{SE} = 0.030$, $t = 1.15$) under Random Forest and
$+0.081$ ($\mathrm{SE} = 0.031$, $p < 0.01$) under XGBoost, indicating
that pilot cities experience measurably faster vegetation recovery
after extreme weather. This finding is consistent with the policy's
emphasis on ecological restoration infrastructure---such as
green corridors, riparian buffers, and adaptive planting
schemes---which operates on post-disturbance recovery rather than
on real-time buffering of NDVI during the event itself.

\subsection{DML vs.\ TWFE-DID Comparison}

To assess the value added by the DML approach, we compare the DML
estimates with conventional TWFE-DID estimates. The two methods produce
similar point estimates (CSEE: DML $= 0.015$ vs.\ TWFE $= 0.017$;
RC: DML $= 0.035$ vs.\ TWFE $= 0.037$), but the DML standard errors are
approximately 25\% smaller, reflecting the efficiency gains from
flexible nuisance function estimation. This comparison suggests that
the DML approach provides meaningful efficiency improvements without
introducing systematic bias, supporting its use as the primary
estimator.
